\documentclass[12pt,a4paper]{article}
\usepackage[utf8]{inputenc}
\usepackage[T5]{fontenc}
\usepackage[vietnamese]{babel}
\usepackage{amsmath}
\usepackage{algorithm}
\usepackage{algpseudocode}
\usepackage{amsfonts}
\usepackage{enumitem}
\usepackage{graphicx}
\begin{document}

\section*{Example 7}
\begin{verbatim}
1 public void function(int n){
2 int i, j, k, count = 0;
3 for(i = n/2; i <= n; i++)
4 for(j = 1; j + n/2 <= n; j++)
5 for(k = 1; k <= n; k = k * 2)
6 count++;
7 }
\end{verbatim}

\textbf{Complexity analysis:}
\begin{itemize}
    \item Loop i: chạy từ $n/2$ đến $n$ $\Rightarrow$ khoảng $n/2$ lần $\Rightarrow O(n)$.
    \item Loop j: điều kiện $j + n/2 \le n \Rightarrow j \le n/2$, $j$ bắt đầu từ 1, mỗi lần tăng 1 $\Rightarrow O(n/2) = O(n)$.
    \item Loop k: $k$ chạy 1, 2, 4, ..., $n$ $\Rightarrow O(\log n)$.
    \item Tổng độ phức tạp: $O(n \times n \times \log n) = O(n^2 \log n)$.
\end{itemize}

\section*{Example 8}
\begin{verbatim}
1 public void function(int n) {
2 int i, j, k, count = 0;
3 for(i = n/2; i <= n; i++)
4 for(j = 1; j <= n; j = 2 * j)
5 for(k = 1; k <= n; k = k * 2)
6 count++;
7 }
\end{verbatim}

\textbf{Complexity analysis:}
\begin{itemize}
    \item Loop i: $O(n)$.
    \item Loop j: $j$ nhân đôi mỗi bước $\Rightarrow O(\log n)$.
    \item Loop k: $O(\log n)$.
    \item Tổng: $O(n \times \log n \times \log n) = O(n \log^2 n)$.
\end{itemize}

\section*{Example 9}
Xác định độ phức tạp của các hàm sau:

\begin{itemize}
    \item[(a)] $T(n) = n\log n + 3n + 2$ $\Rightarrow$ Bậc cao nhất $n\log n$ $\Rightarrow O(n\log n)$.
    \item[(b)] $T(n) = n\log(n!) + 5n^2 + 6$ $\Rightarrow$ $n\log(n!) \approx n(n\log n) = n^2\log n$ (theo Stirling), bậc cao nhất $n^2\log n$ $\Rightarrow O(n^2\log n)$.
    \item[(c)] $T(n) = 100n + 10n^2$ $\Rightarrow$ Bậc cao nhất $n^2$ $\Rightarrow O(n^2)$.
    \item[(d)] $T(n) = 50n\log n + n^3 + 100n$ $\Rightarrow$ Bậc cao nhất $n^3$ $\Rightarrow O(n^3)$.
    \item[(e)] $T(n) = 10n\log n + n(\log n)^2$ $\Rightarrow$ Bậc cao nhất $n(\log n)^2$ $\Rightarrow O(n \log^2 n)$.
\end{itemize}

\section*{Example 10}

\subsection*{a)}
\begin{verbatim}
1 public static int example1(int[] arr) {
2 int n = arr.length, total = 0;
3 for (int j=0; j < n; j++)
4 total += arr[j];
5 return total;
6 }
\end{verbatim}
Độ phức tạp: $O(n)$.

\subsection*{b)}
\begin{verbatim}
1 public static int example2(int[] arr) {
2 int n = arr.length, total = 0;
3 for (int j=0; j < n; j += 2)
4 total += arr[j];
5 return total;
6 }
\end{verbatim}
Vòng lặp chạy $n/2$ lần $\Rightarrow O(n)$.

\subsection*{c)}
\begin{verbatim}
1 public static int example3(int[] arr) {
2 int n = arr.length, total = 0;
3 for (int j=0; j < n; j++)
4 for (int k=0; k <= j; k++)
5 total += arr[j];
6 return total;
7 }
\end{verbatim}
Tổng số lần lặp: $\sum_{j=0}^{n-1} (j+1) = \frac{n(n+1)}{2} \Rightarrow O(n^2)$.

\subsection*{d)}
\begin{verbatim}
1 public static int example4(int[] arr) {
2 int n = arr.length, prefix = 0, total = 0;
3 for (int j=0; j < n; j++) {
4 prefix += arr[j];
5 total += prefix;
6 }
7 return total;
8 }
\end{verbatim}
Mỗi vòng lặp $O(1)$, tổng $n$ lần $\Rightarrow O(n)$.

\subsection*{e)}
\begin{verbatim}
1 public static int example5(int[] first, int[] second) {
2 int n = first.length, count = 0;
3 for (int i=0; i < n; i++) {
4 int total = 0;
5 for (int j=0; j < n; j++)
6 for (int k=0; k <= j; k++)
7 total += first[k];
8 if (second[i] == total) count++;
9 }
10 return count;
11 }
\end{verbatim}
\begin{itemize}
    \item Vòng lặp i: $n$ lần.
    \item Vòng lặp j: $n$ lần.
    \item Vòng lặp k: trung bình $O(n)$ lần (chạy từ 0 đến j).
    \item Mỗi lần lặp trong cùng làm $O(1)$ công việc.
    \item Tổng: $O(n \times n \times n) = O(n^3)$.
\end{itemize}


\section*{Câu 4}

\subsection*{a. Chọn găng tay}
Có tất cả:
\begin{itemize}
    \item 5 đôi màu đỏ $\Rightarrow$ 10 chiếc đỏ.
    \item 4 đôi màu vàng $\Rightarrow$ 8 chiếc vàng.
    \item 2 đôi màu xanh lá cây $\Rightarrow$ 4 chiếc xanh.
\end{itemize}
Tổng số găng tay: $10+8+4 = 22$ chiếc.

\paragraph{Trường hợp tốt nhất:}
Chỉ cần chọn 2 chiếc ngay từ đầu đã tạo thành một đôi.
\[
\boxed{2}
\]

\paragraph{Trường hợp xấu nhất:}
Ta cần xét số lượng găng tay tối đa có thể chọn mà không có đôi nào hoàn chỉnh.  
Mỗi đôi gồm 2 chiếc giống nhau. Để không có đôi nào, ta chỉ lấy tối đa 1 chiếc từ mỗi đôi.  
Có tổng số đôi: $5+4+2 = 11$ đôi.  
Số găng tay tối đa không có đôi là $11$ (mỗi đôi lấy 1 chiếc).  
Thêm 1 chiếc nữa (chiếc thứ 12) sẽ bắt buộc tạo thành một đôi với một trong 11 chiếc đã lấy.  

Vậy trường hợp xấu nhất cần chọn:
\[
\boxed{12}
\]

\subsection*{b. Tất bị mất}
Có 5 đôi tất khác nhau $\Rightarrow$ 10 chiếc tất.  
Mất 2 chiếc tất bất kỳ, mỗi chiếc có xác suất mất như nhau.  
Số cách chọn 2 chiếc tất bị mất từ 10 chiếc:
\[
\binom{10}{2} = 45
\]

\paragraph{Kịch bản tốt nhất:} Mất 2 chiếc cùng một đôi $\Rightarrow$ còn lại 4 đôi hoàn chỉnh.  
Số cách chọn 2 chiếc cùng một đôi: chọn 1 trong 5 đôi $\Rightarrow$ 5 cách.  
Xác suất kịch bản tốt nhất:
\[
P_{\text{tốt nhất}} = \frac{5}{45} = \frac{1}{9} \approx 0.1111
\]

\paragraph{Kịch bản xấu nhất:} Mất 2 chiếc ở 2 đôi khác nhau, mỗi đôi mất 1 chiếc $\Rightarrow$ còn lại 3 đôi hoàn chỉnh và 2 đôi mất mỗi đôi 1 chiếc.  
Số cách: chọn 2 đôi khác nhau từ 5 đôi $\Rightarrow \binom{5}{2} = 10$, mỗi đôi mất 1 trong 2 chiếc $\Rightarrow 2 \times 2 = 4$ cách cho hai đôi đó.  
Tổng số cách:
\[
10 \times 4 = 40
\]
Xác suất kịch bản xấu nhất:
\[
P_{\text{xấu nhất}} = \frac{40}{45} = \frac{8}{9} \approx 0.8889
\]

\paragraph{Số đôi tất kỳ vọng còn lại:}
Gọi $X$ là số đôi còn lại.  
$X$ có thể là 4 (tốt nhất) hoặc 3 (xấu nhất).  
\[
E[X] = 4 \times \frac{1}{9} + 3 \times \frac{8}{9} = \frac{4 + 24}{9} = \frac{28}{9} \approx 3.1111
\]
\[
\boxed{E[X] = \frac{28}{9}}
\]

\section*{Câu 10: Sự tích bàn cờ}

\subsection*{a. Gấp đôi số hạt mỗi ô}
Số hạt ở ô thứ $k$ (bắt đầu ô 1): $2^{k-1}$.  
Tổng số hạt trên 64 ô:
\[
S = 2^0 + 2^1 + \dots + 2^{63} = 2^{64} - 1
\]
Mỗi hạt đếm mất 1 giây.  
Thời gian đếm:
\[
T = 2^{64} - 1 \text{ giây}
\]
Quy đổi:  
1 năm = $365 \times 24 \times 3600 = 31,536,000$ giây.  
\[
T \approx \frac{2^{64}}{3.1536 \times 10^7} \text{ năm}
\]
$2^{64} \approx 1.84467 \times 10^{19}$.  
\[
T \approx \frac{1.84467 \times 10^{19}}{3.1536 \times 10^7} \approx 5.85 \times 10^{11} \text{ năm}
\]
\[
\boxed{\approx 585 \text{ tỷ năm}}
\]

\subsection*{b. Cộng thêm 2 hạt mỗi ô}
Ô thứ nhất: 1 hạt.  
Ô thứ hai: $1+2 = 3$ hạt.  
Ô thứ ba: $3+2 = 5$ hạt.  
Tổng quát: ô thứ $k$ có $1 + 2(k-1) = 2k-1$ hạt.  

Tổng số hạt trên 64 ô:
\[
S = \sum_{k=1}^{64} (2k-1) = 2\sum_{k=1}^{64} k - \sum_{k=1}^{64} 1
\]
\[
= 2 \times \frac{64 \times 65}{2} - 64 = 64 \times 65 - 64 = 64 \times 64 = 4096
\]
Thời gian đếm:
\[
T = 4096 \text{ giây}
\]
Đổi ra phút: $\frac{4096}{60} \approx 68.27$ phút $\approx 1$ giờ 8 phút 16 giây.  
\[
\boxed{4096 \text{ giây (khoảng 1 giờ 8 phút)}}
\]



\section*{Bài tập 1: Kích thước đầu vào và bậc tăng trưởng}

\textbf{Kích thước đầu vào} thường là số lượng phần tử hoặc lượng dữ liệu được truyền vào thuật toán (ví dụ: độ dài mảng, số đỉnh của đồ thị, số bit của một số nguyên).

\textbf{Bậc tăng trưởng} mô tả thời gian chạy hoặc bộ nhớ sử dụng thay đổi thế nào khi kích thước đầu vào tăng lên, bỏ qua các hằng số và các số hạng bậc thấp. Bậc tăng trưởng được biểu diễn bằng các ký hiệu tiệm cận.

\subsection*{Các bậc tăng trưởng thường gặp (từ chậm đến nhanh)}
\[
O(1) \subset O(\log n) \subset O(\sqrt{n}) \subset O(n) \subset O(n\log n) \subset O(n^2) \subset O(2^n) \subset O(n!)
\]

\subsection*{Ví dụ: Xác định kích thước đầu vào và bậc tăng trưởng}
\begin{enumerate}[label=\arabic*.]
    \item \textbf{Tính tổng các phần tử mảng:}
    \begin{verbatim}
    int tong = 0;
    for (int i = 0; i < n; i++) tong += arr[i];
    \end{verbatim}
    Kích thước đầu vào: $n$ (độ dài mảng).  
    Bậc tăng trưởng: $O(n)$ (tuyến tính).

    \item \textbf{Tìm kiếm nhị phân:}
    \begin{verbatim}
    while (left <= right) {
        mid = (left + right) / 2;
        if (arr[mid] == x) return mid;
        else if (arr[mid] < x) left = mid + 1;
        else right = mid - 1;
    }
    \end{verbatim}
    Kích thước đầu vào: $n$ (độ dài mảng).  
    Bậc tăng trưởng: $O(\log n)$ (logarit).

    \item \textbf{Sắp xếp nổi bọt (Bubble sort):}
    \begin{verbatim}
    for (int i = 0; i < n-1; i++)
        for (int j = 0; j < n-i-1; j++)
            if (arr[j] > arr[j+1]) swap(arr[j], arr[j+1]);
    \end{verbatim}
    Kích thước đầu vào: $n$ (độ dài mảng).  
    Bậc tăng trưởng: $O(n^2)$ (bậc hai).

    \item \textbf{Tháp Hà Nội (đệ quy):}
    \begin{verbatim}
    void thapHaNoi(int n, char a, char b, char c) {
        if (n == 1) return;
        thapHaNoi(n-1, a, c, b);
        thapHaNoi(n-1, c, b, a);
    }
    \end{verbatim}
    Kích thước đầu vào: $n$ (số đĩa).  
    Bậc tăng trưởng: $O(2^n)$ (hàm mũ).
\end{enumerate}

\section*{Bài tập 2: Mối quan hệ giữa các độ phức tạp dùng ký hiệu tiệm cận}

Các ký hiệu tiệm cận ($O$, $\Omega$, $\Theta$, $o$, $\omega$) dùng để mô tả mối quan hệ giữa các hàm biểu diễn độ phức tạp thuật toán.

\subsection*{Định nghĩa}
\begin{itemize}
    \item $f(n) = O(g(n))$: Tồn tại $c > 0, n_0 > 0$ sao cho $0 \le f(n) \le c\cdot g(n)$ với mọi $n \ge n_0$. (Chặn trên)
    \item $f(n) = \Omega(g(n))$: Tồn tại $c > 0, n_0 > 0$ sao cho $0 \le c\cdot g(n) \le f(n)$ với mọi $n \ge n_0$. (Chặn dưới)
    \item $f(n) = \Theta(g(n))$: $f(n) = O(g(n))$ và $f(n) = \Omega(g(n))$. (Chặn chặt)
    \item $f(n) = o(g(n))$: $\lim_{n \to \infty} f(n)/g(n) = 0$. (Chặn trên ngặt)
    \item $f(n) = \omega(g(n))$: $\lim_{n \to \infty} f(n)/g(n) = \infty$. (Chặn dưới ngặt)
\end{itemize}

\subsection*{Mối quan hệ giữa các độ phức tạp thường gặp}
Với $n \to \infty$, ta có:

\[
1 = o(\log n) = o(\sqrt{n}) = o(n) = o(n\log n) = o(n^2) = o(2^n)
\]
\[
\log n = o(\sqrt{n}), \quad \sqrt{n} = o(n), \quad n = o(n\log n), \quad n\log n = o(n^2), \quad n^2 = o(2^n)
\]

\subsection*{Ví dụ về mối quan hệ}
\begin{enumerate}[label=\arabic*.]
    \item $5n + 3 = \Theta(n)$
    \item $n^2 + 100n = O(n^3)$ và đồng thời $n^2 + 100n = \Omega(n^2)$
    \item $n\log n = O(n^2)$ nhưng $n\log n \neq \Theta(n^2)$
    \item $2^{n+1} = \Theta(2^n)$
    \item $\log(n!) = \Theta(n\log n)$ (theo xấp xỉ Stirling)
    \item $n^{0.001}$ và $\log n$: thực tế $\log n$ tăng chậm hơn $n^{0.001}$ nên:
    \[
    \lim_{n\to\infty} \frac{\log n}{n^{0.001}} = 0 \quad \Rightarrow \quad \log n = o(n^{0.001})
    \]
    Suy ra $n^{0.001} = \omega(\log n)$.
\end{enumerate}

\subsection*{Bảng tóm tắt mối quan hệ (đúng với $n$ lớn)}
\[
\begin{array}{c|c|c|c|c|c}
f(n) & g(n) & f=O(g) & f=\Omega(g) & f=\Theta(g) & f=o(g) \\
\hline
1000 & 1 & \text{có} & \text{không} & \text{không} & \text{không} \\
\log n & n^{0.1} & \text{có} & \text{không} & \text{không} & \text{có} \\
n & n\log n & \text{có} & \text{không} & \text{không} & \text{có} \\
n\log n & n^2 & \text{có} & \text{không} & \text{không} & \text{có} \\
n^2 & n\log n & \text{không} & \text{có} & \text{không} & \text{không} \\
2^n & n^{100} & \text{không} & \text{có} & \text{không} & \text{không} \\
\end{array}
\]

\section*{Câu 11: Nhẹ hơn hay nặng hơn?}

\textbf{Bài toán:} Có $n > 2$ đồng xu, 1 đồng xu giả (không biết nhẹ hơn hay nặng hơn), biết đồng xu nào là giả. Cần xác định đồng xu giả nhẹ hơn hay nặng hơn với số lần cân $\Theta(1)$.

\textbf{Thuật toán:}
\begin{enumerate}
    \item Lấy đồng xu giả $G$ và 2 đồng xu thật $T_1, T_2$ (lấy từ $n-1$ đồng còn lại, vì $n>2$ nên có ít nhất 2 đồng thật).
    \item Cân $G$ với $T_1$:
    \begin{itemize}
        \item Nếu $G = T_1$: $G$ là thật $\Rightarrow$ vô lý (đã biết $G$ là giả) $\Rightarrow$ không xảy ra.
        \item Nếu $G < T_1$: $G$ nhẹ hơn.
        \item Nếu $G > T_1$: $G$ nặng hơn.
    \end{itemize}
\end{enumerate}
Chỉ cần \textbf{1 lần cân} $\Rightarrow$ $\Theta(1)$.

\begin{algorithm}[H]
\caption{Xác định đồng xu giả nhẹ hay nặng}
\begin{algorithmic}
\REQUIRE $n>2$ đồng xu, biết đồng xu giả $G$, các đồng còn lại là thật
\ENSURE Kết luận "nhẹ hơn" hoặc "nặng hơn"
\STATE Lấy hai đồng thật $T_1, T_2$ từ các đồng còn lại
\STATE Cân $G$ với $T_1$
\IF{$G < T_1$}
    \RETURN "Đồng xu giả nhẹ hơn"
\ELSIF{$G > T_1$}
    \RETURN "Đồng xu giả nặng hơn"
\ENDIF
\end{algorithmic}
\end{algorithm}

\section*{Câu 12: Cánh cửa trên tường}

\textbf{Bài toán:} Đứng trước tường vô tận, cửa ở khoảng cách $n$ bước (chưa biết), không biết hướng. Cần tìm cửa với số bước đi $O(n)$.

\textbf{Thuật toán:} \textbf{Tìm kiếm theo cấp số nhân (Exponential search / "Alternating doubling")}
\begin{enumerate}
    \item Bắt đầu tại vị trí 0, đặt $k = 1$.
    \item Lặp:
    \begin{itemize}
        \item Đi sang phải $k$ bước, kiểm tra cửa.
        \item Nếu chưa thấy, quay về vị trí 0.
        \item Đi sang trái $k$ bước, kiểm tra cửa.
        \item Nếu chưa thấy, quay về vị trí 0.
        \item Nhân đôi $k$: $k = 2k$.
    \end{itemize}
    \item Dừng khi tìm thấy cửa.
\end{enumerate}

\textbf{Chứng minh độ phức tạp:}
Giả sử cửa cách vị trí ban đầu $n$ bước (theo một hướng nào đó). Khi $k$ đạt đến $2^{\lceil \log_2 n \rceil} \le 2n$, thuật toán sẽ thăm hướng đó. Tổng số bước đi:
\[
\sum_{i=0}^{\lceil \log_2 n \rceil} 2\cdot 2^i \cdot 2 = 4 \cdot (2^{\lceil \log_2 n \rceil+1} - 1) \le 4(4n - 1) = O(n)
\]
\boxed{\text{Tổng số bước } \le 16n \text{ bước (cận trên thực tế)}}

\begin{algorithm}[H]
\caption{Tìm cửa trên tường vô tận}
\begin{algorithmic}
\REQUIRE Vị trí ban đầu 0, không biết hướng và khoảng cách $n$
\ENSURE Đến được cửa
\STATE $k \leftarrow 1$
\LOOP
    \FOR{$i \leftarrow 1$ \TO $k$} \STATE đi phải 1 bước \ENDFOR
    \IF{thấy cửa} \RETURN \ENDIF
    \FOR{$i \leftarrow 1$ \TO $k$} \STATE đi trái 1 bước \ENDFOR
    \IF{thấy cửa} \RETURN \ENDIF
    \FOR{$i \leftarrow 1$ \TO $k$} \STATE đi trái 1 bước \ENDFOR
    \IF{thấy cửa} \RETURN \ENDIF
    \FOR{$i \leftarrow 1$ \TO $k$} \STATE đi phải 1 bước \ENDFOR
    \STATE $k \leftarrow 2k$
\ENDLOOP
\end{algorithmic}
\end{algorithm}

\section*{Bài tập 3: Sắp xếp chọn (Selection Sort)}

\subsection*{Mã giả}
\begin{algorithm}[H]
\caption{Selection Sort}
\begin{algorithmic}
\REQUIRE Mảng $a[1..n]$
\ENSURE Mảng $a$ được sắp xếp tăng dần
\FOR{$i \leftarrow 1$ \TO $n-1$}
    \STATE $minIndex \leftarrow i$
    \FOR{$j \leftarrow i+1$ \TO $n$}
        \IF{$a[j] < a[minIndex]$}
            \STATE $minIndex \leftarrow j$
        \ENDIF
    \ENDFOR
    \STATE Hoán đổi $a[i]$ và $a[minIndex]$
\ENDFOR
\end{algorithmic}
\end{algorithm}

\subsection*{Tại sao chỉ cần $n-1$ bước?}
Sau $n-1$ bước, $n-1$ phần tử đầu đã đúng vị trí. Phần tử cuối cùng tự động đúng vì chỉ còn một vị trí trống.

\subsection*{Độ phức tạp}
\begin{itemize}
    \item Best case: $O(n^2)$ (vẫn phải tìm min ở mỗi bước)
    \item Worst case: $O(n^2)$
    \item Trung bình: $O(n^2)$
\end{itemize}
\textbf{So sánh:} Mọi trường hợp đều $O(n^2)$ vì luôn duyệt toàn bộ phần còn lại để tìm min.

\subsection*{Chứng minh tính đúng đắn bằng bất biến lặp}
\textbf{Bất biến lặp:} Sau mỗi lần lặp $i$, đoạn $a[1..i]$ chứa $i$ phần tử nhỏ nhất của mảng ban đầu và được sắp xếp tăng dần.

\begin{itemize}
    \item \textbf{Khởi tạo:} $i=1$, $a[1..1]$ hiển nhiên đúng.
    \item \textbf{Duy trì:} Giả sử đúng sau $i-1$. Ở bước $i$, tìm $min$ trong $a[i..n]$ và đưa về $a[i]$. Khi đó $a[1..i]$ gồm $i$ phần tử nhỏ nhất và sắp xếp.
    \item \textbf{Kết thúc:} Khi $i=n-1$, $a[1..n-1]$ đúng, phần tử cuối tự động đúng $\Rightarrow$ mảng sắp xếp hoàn chỉnh.
\end{itemize}

\subsection*{Cài đặt và kiểm thử (mã Python)}
\begin{verbatim}
import random
import time
import matplotlib.pyplot as plt

def selection_sort(arr):
    n = len(arr)
    for i in range(n-1):
        min_idx = i
        for j in range(i+1, n):
            if arr[j] < arr[min_idx]:
                min_idx = j
        arr[i], arr[min_idx] = arr[min_idx], arr[i]
    return arr

# Đo thời gian
sizes = [100, 200, 500, 1000, 2000, 5000]
times = []
for size in sizes:
    arr = [random.randint(0, 10000) for _ in range(size)]
    start = time.time()
    selection_sort(arr.copy())
    end = time.time()
    times.append(end - start)

# Vẽ biểu đồ
plt.plot(sizes, times, 'o-')
plt.xlabel('Kích thước đầu vào (n)')
plt.ylabel('Thời gian (giây)')
plt.title('Selection Sort - Thời gian chạy')
plt.grid(True)
plt.show()
\end{verbatim}

\begin{center}
\includegraphics[width=0.7\textwidth]{selection_sort_plot.png}
\end{center}

\subsection*{Đề xuất hai thuật toán sắp xếp khác}

\subsubsection*{1. Sắp xếp chèn (Insertion Sort)}
\begin{algorithm}[H]
\caption{Insertion Sort}
\begin{algorithmic}
\FOR{$i \leftarrow 2$ \TO $n$}
    \STATE $key \leftarrow a[i]$
    \STATE $j \leftarrow i-1$
    \WHILE{$j \ge 1$ AND $a[j] > key$}
        \STATE $a[j+1] \leftarrow a[j]$
        \STATE $j \leftarrow j-1$
    \ENDWHILE
    \STATE $a[j+1] \leftarrow key$
\ENDFOR
\end{algorithmic}
\end{algorithm}
Độ phức tạp: Best $O(n)$, Worst $O(n^2)$.

\subsubsection*{2. Sắp xếp nhanh (Quick Sort)}
\begin{algorithm}[H]
\caption{Quick Sort}
\begin{algorithmic}
\PROCEDURE{QuickSort}{$a, low, high$}
    \IF{$low < high$}
        \STATE $pivotIndex \leftarrow \text{Partition}(a, low, high)$
        \STATE \CALL{QuickSort}{$a, low, pivotIndex-1$}
        \STATE \CALL{QuickSort}{$a, pivotIndex+1, high$}
    \ENDIF
\ENDPROCEDURE
\end{algorithmic}
\end{algorithm}
Độ phức tạp: Best $O(n\log n)$, Worst $O(n^2)$, trung bình $O(n\log n)$.

\subsection*{So sánh ba thuật toán (đồ thị)}
Sinh viên tự cài đặt và vẽ đồ thị so sánh thời gian chạy giữa Selection Sort, Insertion Sort và Quick Sort trên cùng bộ dữ liệu.

\end{document}